% SOURCE_AUTHORITY: sga5_fr_workpass.tex, lines 12971--13016 (LF-normalized unit).
% SOURCE_SHA256: 791F4EFFC5E02832D5D77ED03518C8156D6F07E4C8238B03545DB93D883FBB28
c) Seja $\Phi$ um complexo de feixes como em b). Tem-se o isomorfismo
\[
\RG_C(\Phi)=(\R\Gamma^G)(\RG_{C'}(\Phi')),
\qquad \Phi'=p^*(\Phi).
\tag{7.5}
\]
Com efeito:
\[
(\R\Gamma^G)(\RG_{C'}(\Phi'))
\simeq
(\R\Gamma^G)(\RG_C(p_*(\Phi')))
\simeq
\RG_C(\R\underline{\Gamma}^G(p_*(\Phi')))
\]
graças a (7.3), e (7.5) resulta do isomorfismo (7.4).

d) Denotemos por $i':U'\to C'$, $i:U\to C$ as inclusões canônicas, e ponhamos
\[
\Phi=F_{U,C}=i_!(F|U).
\]
Para calcular $\RG_C(\Phi)$ pode-se aplicar a fórmula (7.5). Por outro lado, com as notações de 7.1,
\[
\Phi'\simeq \Lambda_{U',C'}\otimes_\Lambda^L M,
\]
e a fórmula de Künneth (SGA 4 XVII) dá
\[
\RG_{C'}(\Phi')
\simeq
\RG_{C'}(\Lambda_{U',C'}\otimes_\Lambda^L M)
\xleftarrow{\sim}
\RG_{C'}(\Lambda_{U',C'})\otimes_\Lambda^L M.
\tag{7.6}
\]
Este último isomorfismo é, de fato, um isomorfismo entre os objetos de $D(A[G])$ definidos por ambos os membros. Combinando (7.6), (7.5) e a fórmula (3.9.1), que é aplicável graças a 2.1, obtém-se
\[
\RG_C(F_{U,C})
\simeq
\R\Gamma^G\bigl(\RG_{C'}(\Lambda_{U',C'})\otimes_\Lambda^L M\bigr)
\simeq
\RHom_{\Lambda[G]}\bigl(\RG_{C'}(\Lambda_{U',C'})^\vee,M\bigr).
\tag{7.7}
\]
Por conseguinte, tendo em conta que $M\simeq F_{\bar\eta}\in\Ob(\Delta)$ (condição (i)) e (3.6), obtém-se
\[
\RG_C(F_{U,C})\in\Ob(\Delta).
\]
